\subsection{Almost complex manifolds and complex manifolds}

In this section, we assume K = R and denote by X a (real) manifold of class \(\mathbf{C}^{r}\) , with \(r \in \mathbf{N}_{\mathbf{R}}\) .

8.8.1 (“Complex vector bundles”). A vector bundle over \(\mathbf{C}\) with base X (7.3.4) is also called a complex vector bundle with base X. If M is such a vector bundle, one calls a complex vector chart of M a triple (U, \(\varphi\) , E), where U is an open subset of X, E a complex Banach space, and \(\varphi\) an isomorphism of complex vector bundles from \(M|_U\) onto the trivial bundle \(E_U = U \times E\). A family \((U_i,\varphi_i,E_i)\) of complex vector charts on M whose domains \(U_i\) cover X is called a complex vector atlas of M; every complex vector bundle possesses a complex vector atlas (7.3.3).

One sometimes says C-vector instead of complex vector. Similarly, when one wishes to speak of a vector bundle over R, of an ordinary vector chart of this bundle, etc., one sometimes says "real vector" or "R-vector" instead of "vector".

Let M be a complex vector bundle with base X. There exists one and only one automorphism j of the (real) vector bundle underlying M (7.3.3) whose restriction to each fibre is multiplication by the element i of C; the square of j is equal to -1. Conversely, if M is a (real) vector bundle with base X and j is an automorphism of M such that \(j^{2} = -1\) , there exists on M a unique structure of complex vector bundle admitting M as underlying (real) vector bundle and for which j is multiplication by i.

8.8.2 ("Complexification of a vector bundle"). One defines a vector functor \(\tau\) by setting \(\tau(E) = E \otimes C\) for every (real) Banach space \(E\) , and \(\tau(u) = u \otimes Id_{C}\) for every morphism \(u: E \to E'\) of (real) Banach spaces.

If F is a vector bundle with base X and of class \(C^{k}\) ( \(0 \leqslant k \leqslant r\) ), its image under the vector functor \(\tau\) is denoted \(F \otimes C\) . There exists on \(F \otimes C\) a unique structure of complex vector bundle with base X and of class \(C^{k}\) , compatible with its structure of (real) vector bundle and inducing on each fibre \((\mathbf{F} \otimes \mathbf{C})_{x} = \mathbf{F}_{x} \otimes \mathbf{C}\) its canonical structure of complex Banach space (EVT, II, § 8, n° 1, Example). Endowed with this structure, one says that \(F \otimes C\) is the complexification of F.

Let U be an open subset of X; the canonical map

\[
\mathcal {S} _ {\mathrm{F}} ^ {k} (\mathrm{U}) \otimes \mathbf {C} \rightarrow \mathcal {S} _ {\mathrm{F} \otimes \mathbf {C}} ^ {k} (\mathrm{U})
\]

is an isomorphism, by which one identifies these two spaces. If \(s \in \mathcal{S}_{\mathrm{F} \otimes \mathrm{C}}^{k}(\mathrm{U})\) , the elements \(\sigma, \tau\) of \(\mathcal{S}_{\mathrm{F}}^{k}(\mathrm{U})\) such that \(s = \sigma + i\tau\) are called the real part and the imaginary part of \(s\) ; the section \(\bar{s} = \sigma - i\tau\) is called the conjugate of \(s\) .

In particular, a section of the bundle \(\mathrm{T}(\mathrm{X}) \otimes \mathbf{C}\) is called a complex vector field on X.

Let H be a complex vector bundle over the base X. The canonical isomorphisms \(\operatorname{Alt}_{\mathbf{R}}^{p}(\mathrm{T}_{x}(\mathrm{X});\mathrm{H})\to\operatorname{Alt}_{\mathbf{C}}^{p}(\mathrm{T}_{x}(\mathrm{X})\otimes\mathbf{C};\mathbf{H})\) define a canonical isomorphism of complex vector bundles, called the canonical one, from \(\operatorname{Alt}_{\mathbf{R}}^{p}(\mathrm{T}(\mathrm{X});\mathrm{H})\) onto \(\operatorname{Alt}_{\mathbf{C}}^{p}(\mathrm{T}(\mathrm{X})\otimes\mathbf{C};\mathbf{H})\) (7.8.5), by means of which we will identify these two bundles. Let \(\omega\) be a section of this bundle, in other words a differential form of degree p on X, with values in H, and let \(\zeta\) be a complex vector field, with real part \(\xi\) and imaginary part \(\eta\). We then set:

\[
i (\zeta , \omega) = i (\xi , \omega) + i. i (\eta , \omega).
\]

If H is a trivial bundle, we set

\[
\theta_ {\zeta}. \omega = \theta_ {\xi}. \omega + i \theta_ {\eta}. \omega .
\]

Similarly, we extend the bracket by linearity to complex vector fields. The formulas of the preceding numbers remain valid.

8.8.3. An almost complex structure of class \(C^{k}\) ( \(k \leqslant r - 1\) ) on X is the data of a complex vector bundle structure on \(\mathrm{T}(\mathrm{X})\), of class \(C^{k}\) , whose underlying real structure is the natural structure of \(\mathrm{T}(\mathrm{X})\). Such a structure is equivalent to the data of an automorphism \(j\) of class \(C^{k}\) of \(\mathrm{T}(\mathrm{X})\) such that \(j^{2} = -1\) .

Let us take such a structure. Extend \(j\) to a \(\mathbf{C}\)-automorphism of \(\mathrm{T}(\mathrm{X}) \otimes \mathbf{C}\). There exists a unique decomposition of \(\mathrm{T}(\mathrm{X}) \otimes \mathbf{C}\) as a direct sum of two complex subbundles:

\[
\mathbf {T} (\mathrm{X}) \otimes \mathbf {C} = \mathbf {T} ^ {\prime} (\mathrm{X}) \oplus \mathbf {T} ^ {\prime \prime} (\mathrm{X})\tag{1}
\]

such that \(j\) coincides with multiplication by \(i\) on \(\mathbf{T}'(\mathbf{X})\) and by \(-i\) on \(\mathbf{T}''(\mathbf{X})\). The projector \(p'\) of \(\mathbf{T}(\mathbf{X}) \otimes \mathbf{C}\) onto \(\mathbf{T}'(\mathbf{X})\) is equal to \(\frac{1}{2}(1 - ij)\) ; the projector \(p'' = 1 - p'\) of \(\mathbf{T}(\mathbf{X}) \otimes \mathbf{C}\) onto \(\mathbf{T}''(\mathbf{X})\) is equal to \(\frac{1}{2}(1 + ij)\) . The composites

\[
\pi^ {\prime} \colon \mathrm{T} (\mathrm{X}) \rightarrow \mathrm{T} (\mathrm{X}) \otimes \mathbf {C} \xrightarrow {p ^ {\prime}} \mathrm{T} ^ {\prime} (\mathrm{X})
\]

\[
\pi^ {\prime \prime}: \mathrm{T} (\mathrm{X}) \rightarrow \mathrm{T} (\mathrm{X}) \otimes \mathbf {C} \xrightarrow {p ^ {\prime \prime}} \mathrm{T} ^ {\prime \prime} (\mathrm{X})
\]

are R-isomorphisms; the first maps j to multiplication by i, the second to multiplication by -i.

8.8.4 (“Forms of type \((p,q)\)"). We keep the hypotheses and notations of the preceding number. Let F be a complex vector bundle with base X, and let p and q be two integers \(\geqslant0\) , and let \(n=p+q\) . Let \(\omega\) be a differential form of degree n on an open set U of X, with values in F; identify (8.8.2) \(\omega\) with a section of \(\operatorname{Alt}_{\mathbf{C}}^{n}(\operatorname{T}(\mathbf{X})\otimes\mathbf{C};\mathbf{F})\) . We say that \(\omega\) is of type \((p,q)\) if, for all \(x\in U\) , one has

\[
\omega_ {x} (v _ {1}, \dots , v _ {n}) = 0
\]

as soon as \(p + 1\) vectors \(v_i\) belong to \(T_x'(X)\) or \(q + 1\) vectors \(v_i\) belong to \(T_x''(X)\). If one identifies \(\bigwedge^n T_x(X) \otimes C\) with the direct sum of the spaces \(\bigwedge^{m'} T_x'(X) \otimes \bigwedge^{m''} T_x''(X)\) for \(m' + m'' = n\) (A, III, p. 85) and \(\omega_x\) with a \(C\)-linear form on \(\bigwedge^n T_x(X)\) , it is equivalent to say that \(\omega_x\) vanishes on \(\bigwedge^{m'} T_x'(X) \otimes \bigwedge^{m''} T_x''(X)\) for \((m', m'') \neq (p, q)\) .

For every pair \((p,q)\) of integers \(\geqslant0\) with \(p+q=n\) there exists a complex vector subbundle \(\operatorname{Alt}^{p,q}(\mathrm{T}(\mathrm{X});\mathrm{F})\) of \(\operatorname{Alt}^{n}(\mathrm{T}(\mathrm{X});\mathrm{F})\) and only one such that its sections over an open set U of X are the differential forms of type \((p,q)\) on X, with values in F. The vector bundle \(\operatorname{Alt}^{n}(\mathrm{T}(\mathrm{X});\mathrm{F})\) is the direct sum of the vector bundles \(\operatorname{Alt}^{p,q}(\mathrm{T}(\mathrm{X});\mathrm{F})\) for \(p\geqslant0\) , \(q\geqslant0\) and \(p+q=n\) . Every differential form \(\omega\) of degree n with values in F decomposes uniquely into

\[
\omega = \sum_ {p + q = n} \omega_ {p, q}
\]

where \(\omega_{p,q}\) is a form of type \((p,q)\) , called the component of type \((p,q)\) of \(\omega\) . If \(\omega\) is of class \(\mathbf{C}^h\) , with \(0 \leqslant h \leqslant k\) , the same holds for its components.

Examples

\begin{enumerate}
\def\labelenumi{\alph{enumi})}
\item
  A differential form of degree n is of type \((n, 0)\) if and only if it is \(\mathbf{C}\)-multilinear.
\item
  Let \(\mathbf{F}'\), \(\mathbf{F}''\), and \(\mathbf{F}\) be three complex vector bundles and let \(\mathbf{F}' \times_{\mathbf{X}} \mathbf{F}'' \to \mathbf{F}\) be a \(\mathbf{C}\)-bilinear pairing. If \(\omega'\) (resp. \(\omega''\)) is a differential form of type \((p',q')\) (resp. \((p'',q'')\)) with values in \(\mathbf{F}'\) (resp. \(\mathbf{F}''\)), the differential form \(\omega' \wedge \omega''\) is of type \((p' + p'',q' + q'')\).
\item
  Suppose that F is the complexification of a (real) vector bundle (8.8.2). If \(\omega\) is a differential form of type \((p, q)\) with values in F, its conjugate \(\overline{\omega}\) (8.8.2) is of type \((q, p)\) .
\end{enumerate}

8.8.5. In addition to the preceding hypotheses, we assume \(k \geqslant 1\) . The following three conditions are then equivalent:

\begin{enumerate}
\def\labelenumi{(\roman{enumi})}
\item
  For every open set U of X and for every pair \((\xi, \eta)\) of complex vector fields on U, of class \(C^{k}\) such that \(\xi(x)\) and \(\eta(x)\) belong to \(\mathrm{T}_{x}^{\prime}(\mathrm{X})\) for all \(x \in U\) , one has \([\xi, \eta](x) \in \mathrm{T}_{x}^{\prime}(\mathrm{X})\) for all \(x \in U\) .
\item
  For every open set U of X and every pair \((\xi, \eta)\) of vector fields (real) on U, of class \(C^{k}\) , one has
\end{enumerate}

\[
[ \xi , \eta ] + j [ j \xi , \eta ] + j [ \xi , j \eta ] - [ j \xi , j \eta ] = 0.
\]

\begin{enumerate}
\def\labelenumi{(\roman{enumi})}
\setcounter{enumi}{2}
\tightlist
\item
  For every open set U of X and for every complex differential form \(\omega\) of type \((1,0)\) on U, of class \(C^{k}\) , the component of type \((0,2)\) of \(d\omega\) is zero.
\end{enumerate}

Suppose these conditions are satisfied. Let \(\omega\) a differential form of type \((p,q)\) on an open set U of X, of class \(\mathbf{C}^h\) with \(1\leqslant h\leqslant k\) with values in a complex Banach space F. We denote by \(d'\omega\) (resp. \(d''\omega\) ) the component of type \((p+1,q)\) (resp. \((p,q+1))\) of \(d\omega\) ; the other components of \(d\omega\) are then zero and one has

\[
d \omega = d ^ {\prime} \omega + d ^ {\prime \prime} \omega .\tag{2}
\]

If \(k\geqslant 2\) , one has

\[
d ^ {\prime} (d ^ {\prime} \omega) = 0, d ^ {\prime \prime} (d ^ {\prime \prime} \omega) = 0, d ^ {\prime} (d ^ {\prime \prime} \omega) + d ^ {\prime \prime} (d ^ {\prime} \omega) = 0.\tag{3}
\]

With the notation of 8.8.4, example \(b\) ), we have

\begin{enumerate}
\def\labelenumi{(\arabic{enumi})}
\setcounter{enumi}{3}
\tightlist
\item
\end{enumerate}

\[
d ^ {\prime} \left(\omega^ {\prime} \wedge \omega^ {\prime \prime}\right) = d ^ {\prime} \left(\omega^ {\prime}\right) \wedge \omega^ {\prime \prime} + (- 1) ^ {p ^ {\prime} + q ^ {\prime}} \omega^ {\prime} \wedge d ^ {\prime} \left(\omega^ {\prime \prime}\right)\tag{4}
\]

\[
d ^ {\prime \prime} (\omega^ {\prime} \wedge \omega^ {\prime \prime}) = d ^ {\prime \prime} (\omega^ {\prime}) \wedge \omega^ {\prime \prime} + (- 1) ^ {p ^ {\prime} + q ^ {\prime}} \omega^ {\prime} \wedge d ^ {\prime \prime} (\omega^ {\prime \prime})\tag{5}
\]

Together with those of 8.8.4, example c), we have

\[
d ^ {\prime} (\overline {{\omega}}) = \overline {{d ^ {\prime \prime} (\omega)}}.\tag{6}
\]

In formulas (4) and (5) (resp. (6)), the differential forms \(\omega'\) and \(\omega''\) (resp. \(\omega\) ) are assumed to be of class \(\mathbf{C}^1\) .

8.8.6 (“Almost complex structure of a complex manifold”). Let \(X^c\) be a complex analytic manifold, admitting X as its underlying real analytic manifold (5.14.2.a)). The bundle \(\mathrm{T}(\mathrm{X})\) is identified with the (real) vector bundle underlying the complex vector bundle \(\mathrm{T}(X^c)\), which endows X with an almost complex structure of class \(C^\omega\) , said to be defined by the given complex analytic manifold structure on X. This almost complex structure satisfies conditions (i) to (iii) of 8.8.5; in particular, the operators \(d'\) and \(d''\) of 8.8.5 are defined. If \(f: X^c \to Y^c\) is a morphism of complex analytic manifolds, \(f^*\) commutes with the operators \(d'\) and \(d''\) .

8.8.7. We retain the hypotheses and notation of 8.8.6. Let \(\overline{\mathbf{X}}^c\) be the conjugate of \(\mathbf{X}^c\) (5.14.2.b)). Embed X by the diagonal map \(x \mapsto (x, x)\) into the complex manifold \(\mathbf{X}^c \times \overline{\mathbf{X}}^c\) , and let \(g\) be the anti-automorphism \((x, y) \mapsto (y, x)\) of \(\mathbf{X}^c \times \overline{\mathbf{X}}^c\) . The pair \((\mathbf{X}^c \times \overline{\mathbf{X}}^c, g)\) is a complexification of X (5.14.8). In particular, for every \(x \in \mathbf{X}\) , the tangent space \(\mathrm{T}_x(\mathrm{X}^c \times \overline{\mathrm{X}}^c)\) is identified with \(\mathrm{T}_x(\mathrm{X}) \otimes \mathbf{C}\) ; in this identification, \(\mathrm{T}_x(\mathrm{X}^c \times \{x\})\) corresponds to \(\mathrm{T}_x'(\mathrm{X})\) and \(\mathrm{T}_x(\{x\} \times \overline{\mathrm{X}}^c)\) corresponds to \(\mathrm{T}_x''(\mathrm{X})\) .

8.8.8. Suppose \(r = \omega\) ; any almost complex structure of class \(C^\omega\) on X satisfying the equivalent conditions (i) to (iii) of 8.8.5 is the almost complex structure defined by one and only one complex analytic manifold structure on \(X\) . \(^{(1)}\)

8.8.9 (“Holomorphic forms”). Let us resume the hypotheses and notations of 8.8.6, and let F be a complex Banach space. The morphic differential forms (in other words, complex analytic ones) on \(\mathbf{X}^c\) with values in F are still called holomorphic differential forms; the identification of \(\mathrm{T}(\mathbf{X}^c)\) with \(\mathrm{T}(\mathbf{X})\) identifies them with (real analytic) differential forms on X. For a differential form \(\omega\) of degree \(p\) on X, with values in F and of class \(\mathbf{C}^k\), with \(k \geqslant 1\), for it to be holomorphic, it is necessary and sufficient that it be of type \((p, 0)\) and that \(d''\omega = 0\) .

Let \(\xi\) be a vector field on X. We say that \(\xi\) is holomorphic if it is a complex analytic mapping from \(X^c\) into \(T(X^c) = T(X)\), in other words, a complex analytic vector field on \(X^c\). If this is the case, and if \(\omega\) is a holomorphic differential form on X with values in F, the form \(\theta_{\xi}.\omega\) (resp. \(i(\xi)\omega\) ) is the same, whether one considers \(\xi\) as a complex analytic vector field on \(X^c\) and \(\omega\) as an element of \(\Omega^p(X^c; F)\) , or \(\xi\) as a vector field on X and \(\omega\) as an element of \(\Omega^p(X; F)\) . Moreover, let \(\xi' = \pi'(\xi)\) and \(\xi'' = \pi''(\xi)\) be the components of \(\xi\) in \(T'(X)\) and \(T''(X)\) respectively (8.8.3); we have

\begin{enumerate}
\def\labelenumi{(\arabic{enumi})}
\setcounter{enumi}{6}
\tightlist
\item
\end{enumerate}
\[
\theta_ {\xi}. \omega = \theta_ {\xi^ {\prime}}. \omega , \quad \theta_ {\xi^ {\prime \prime}}. \omega = 0,\tag{7}
\]

\[
i (\xi) \omega = i (\xi^ {\prime}) \omega , i (\xi^ {\prime \prime}) \omega = 0.\tag{8}
\]

8.8.10. Example. --- Suppose that \(X^{c}\) is an open subset of \(C^{n}\) and denote by \(z^{1},\ldots,z^{n}\) the coordinate functions on \(X^{c}\) . For \(1\leqslant k\leqslant n\) , set \(z^{k}=x^{k}+iy^{k}\) , where \(x^{k}\) and \(y^{k}\) take real values. Consider the tangent frame \((\partial/\partial z^{k})\) of \(\mathrm{T}(\mathrm{X}^{c})\) defined by the coordinate system \((z^{k})\) and the tangent frame \((\partial/\partial x^{k},\partial/\partial y^{k})\) of \(\mathrm{T}(\mathrm{X})\) defined by the coordinate system \((x^{k},y^{k})\) (8.1.5). We have \(j(\partial/\partial x^{k})=\partial/\partial y^{k}\) and the image of \(\partial/\partial z^{k}\) under the map \(\pi^{\prime}\colon\mathrm{T}(\mathrm{X})=\mathrm{T}(\mathrm{X}^{c})\to\mathrm{T}^{\prime}(\mathrm{X})\) (8.8.3) is equal to

\[
\frac {1}{2} (\partial / \partial x ^ {k} - i \partial / \partial y ^ {k});
\]

one generally still denotes it by \(\partial/\partial z^{k}\) . \(^{(2)}\) The complex conjugate vector field of this field \(\partial/\partial z^{k}\) is denoted \(\partial/\partial\bar{z}^{k}\) ; we have

\[
\partial / \partial \bar {z} ^ {k} = \frac {1}{2} (\partial / \partial x ^ {k} + i \partial / \partial y ^ {k}).\tag{9}
\]

If f is a function of class \(C^{1}\) on X, with values in a complex Banach space F, the differential forms \(d'f\) and \(d''f(8.8.5)\) are given by

\[
d ^ {\prime} f = \sum_ {1 \leqslant k \leqslant n} \frac {\partial f}{\partial z ^ {k}} d z ^ {k}, d ^ {\prime \prime} f = \sum_ {1 \leqslant k \leqslant n} \frac {\partial f}{\partial \bar {z} ^ {k}} d \bar {z} ^ {k}.\tag{10}
\]

\(^{(1)}\) If X is of finite dimension, this remains true for an almost complex structure of class \(C^k\) with \(k = \dim X\) (cf. A. Newlander and L. Nirenberg, "Complex analytic coordinates in almost complex manifolds," Ann. of Math. LXX (1957), pp. 391--404).

\(^{(2)}\) Care must be taken: the canonical identification of \(\mathrm{T}(X^c)\) with \(\mathrm{T}(X)\) does not transform the vector field \(\partial/\partial z^k\) on \(X^c\) into the vector field \(\partial/\partial z^k\) on X, but into the vector field \(\partial/\partial z^k+\partial/\partial\bar z^k\). However, the two vector fields denoted \(\partial/\partial z^k\) have the same value on holomorphic differential forms (8.8.9), which are in fact the only forms to which one can apply a vector field of \(X^c\).

Recall that we have

\[
d f = d ^ {\prime} f + d ^ {\prime \prime} f
\]

\[
d z ^ {k} = d x ^ {k} + i d y ^ {k} \quad d \bar {z} ^ {k} = d x ^ {k} - i d y ^ {k}.
\]

If \(\mathbf{A}=(i_{1},\ldots,i_{p})\) is a strictly increasing sequence of integers belonging to the interval \([1,n]\) , we set

\[
d z ^ {A} = d z ^ {i _ {1}} \wedge \dots \wedge d z ^ {i _ {p}}
\]

and denote by \(d\bar{z}^{\mathrm{A}}\) the conjugate of \(dz^{\mathrm{A}}\) . Every differential form \(\omega\) of type \((p,q)\) with values in F is written uniquely as

\[
\omega = \sum_ {\mathrm{A}, \mathrm{B}} f _ {\mathrm{A}, \mathrm{B}} d z ^ {\mathrm{A}} \wedge d \bar {z} ^ {\mathrm{B}}\tag{11}
\]

where A (resp. B) ranges over the set of strictly increasing sequences of \(p\) (resp. \(q\) ) elements of \([1, n]\) , the \(f_{A,B}\) being functions with values in F. If \(\omega\) is of class \(\mathbf{C}^k\) , with \(k \geqslant 1\) , the same holds for the functions \(f_{A,B}\) , and one has

\begin{enumerate}
\def\labelenumi{(\arabic{enumi})}
\setcounter{enumi}{11}
\tightlist
\item
\end{enumerate}
\[
d ^ {\prime} \omega = \sum_ {A, B} d ^ {\prime} f _ {A, B} \wedge d z ^ {A} \wedge d \bar {z} ^ {B}\tag{12}
\]

\[
d ^ {\prime \prime} \omega = \sum_ {\mathbf {A}, \mathbf {B}} d ^ {\prime \prime} f _ {\mathbf {A}, \mathbf {B}} \wedge d z ^ {\mathbf {A}} \wedge d \bar {z} ^ {\mathbf {B}}.\tag{13}
\]

